For $x<y\le0$, $x^2>y^2$, while for $0\le x<y$, $x^2<y^2$. Thus $f$ is strictly decreasing on $(-\infty,0]$ and strictly increasing on $[0,\infty)$; it is neither on the whole real line.

The three average rates are $\dfrac{1-4}{1}=-3$, $\dfrac{0-1}{1}=-1$, and $\dfrac{1-0}{1}=1$. More generally,
\[\dfrac{(x+h)^2-x^2}{h}=2x+h,\]
whose limit is $f'(x)=2x$. The derivative is increasing everywhere, so the graph is concave up everywhere. At $x=0$, $f(0)=0$ and $f'(0)=0$, so it is both a zero and a stationary point. It is not an inflection point: the concavity does not change there. A zero derivative alone does not imply an inflection point.
\begin{asy}
texpreamble("\usepackage{amsmath,amssymb}");
size(240);
defaultpen(fontsize(9pt));
import graph;real f(real x){return x*x;}draw((-2.5,0)--(2.5,0),Arrow);draw((0,-.3)--(0,4.8),Arrow);draw(graph(f,-2.1,2.1),blue);draw((-2,4)--(-1,1),red+dashed);draw((-1,1)--(0,0),red+dashed);draw((0,0)--(1,1),red+dashed);label("$x$",(2.5,0),(0,-1));label("$y$",(0,4.8),(1,0));label("$y=x^2$",(1.7,3.5),(1,0));label("$0$",(0,0),(-1,-1));
\end{asy}
